%------------------------------------------------------
%------------------------------------------------------

\section{Detección binaria por familia}
\label{apx:binaria_familia}

La Tabla~\ref{tab:apx_binaria_familia} detalla la exactitud de detección binaria (cifrado vs.\ seguro) por familia obtenida con Random Forest en el Experimento~1, referida en la Sección~\ref{sec:exp1}.

\begin{table}[ht]
\centering
\caption{Exactitud de detección binaria por familia (Random Forest)}
\label{tab:apx_binaria_familia}
\begin{tabular}{|l|c||l|c|}
\hline
\textbf{Familia} & \textbf{Exactitud (\%)} & \textbf{Familia} & \textbf{Exactitud (\%)} \\ \hline
WANNACRY & 97,3 & CERBER & 88,7 \\ \hline
CHIMERA & 96,7 & NOTPETYA & 88,7 \\ \hline
MAZE & 96,7 & CLOP & 88,0 \\ \hline
PHOBOS & 96,7 & AVOSLOCKER & 82,7 \\ \hline
BLACKMATTER & 95,3 & CONTI & 82,7 \\ \hline
DARKSIDE & 95,3 & LOCKBIT & 80,7 \\ \hline
BADRABBIT & 94,7 & BLACKBASTA & 78,0 \\ \hline
DHARMA & 94,7 & CUBA & 78,0 \\ \hline
SODINOKIBI & 94,7 & HELLOKITTY & 76,7 \\ \hline
WASTEDLOCKER & 94,7 & SUNCRYPT & 75,3 \\ \hline
CRYPTOLOCKER & 94,0 & NETWALKER & 74,0 \\ \hline
MEDUZALOCKER & 94,0 & RANSOMEXX & 69,3 \\ \hline
RYUK & 93,3 & & \\ \hline
TESLACRYPT & 93,3 & & \\ \hline
BLACKCAT & 92,0 & & \\ \hline
LORENZ & 91,3 & & \\ \hline
GANDCRAB & 90,7 & & \\ \hline
JIGSAW & 90,7 & & \\ \hline
\end{tabular}
\end{table}

\section{Reproducibilidad}
\label{apx:reproducibilidad}

Cada cifra del Capítulo~4 procede de un archivo de resultados y puede contrastarse con él mediante los dos verificadores que acompañan al código, uno por frente. No son una declaración de principio: leen los números tal como están escritos en el documento y los comparan con su archivo de origen, de modo que una discrepancia aparece sola.

El verificador del frente de archivos cifrados, \texttt{cifras\_finales\_archivos.py}, comprueba 630 cifras y las confirma todas. De ellas, 377 no se leen de un resumen ya calculado sino que se recalculan desde los datos por semilla o por familia, lo que permite detectar un promedio mal tomado y no solo una transcripción equivocada. El verificador del frente de notas, \texttt{cifras\_finales.py}, comprueba 87 y las confirma todas. Ambos emiten además la lista de advertencias necesarias para citar cada número con su métrica y su base.

Los dos declaran de forma explícita las cifras que no provienen de una corrida guardada, junto con su procedencia: la evaluación de herramientas públicas, realizada a mano y registrada en la planilla del proyecto; el censo de integridad, que es un recuento del conjunto de archivos; las pruebas preliminares y la ampliación del Experimento~2, cuyas corridas son anteriores a los verificadores; y la composición de los pliegues de la validación por tipo de documento. Una cifra de esta clase no es menos comprobable, pero se comprueba abriendo su fuente en lugar de ejecutando un guion, y el documento lo señala en cada caso.

Corresponde declarar un límite. Las corridas conservan las métricas agregadas por semilla, no la predicción sobre cada archivo individual, de modo que las cifras que dependerían de esas predicciones se leen del registro de ejecución del experimento y los verificadores las marcan como tales en vez de recalcularlas.

Hay además un rasgo del capítulo que conviene tener presente al recorrerlo. Una misma métrica aparece medida sobre bases distintas, porque el corpus de notas creció y se depuró durante el trabajo y porque varios experimentos son anteriores a esa depuración. Cada cifra se reporta junto a la base sobre la que se midió, y dos cifras de bases distintas no son comparables entre sí ni se corrigen una con otra: una misma métrica puede valer distinto sobre 146 notas y sobre 149 porque son dos mediciones de conjuntos diferentes, y el capítulo lo declara en cada caso. Lo que los verificadores detectan es una transcripción equivocada o una base mal atribuida, no una diferencia legítima entre bases, y por eso una cifra sin su base declarada al lado es el defecto que más interesa encontrar.

Cada experimento escribe un manifiesto con los parámetros del modelo, las semillas empleadas, las versiones de las librerías y, cuando corresponde, el identificador del trabajo en el clúster. El proyecto conserva 51 de esos manifiestos. Los experimentos de mayor costo se ejecutaron en el clúster del NIDTEC mediante SLURM, y sus 21 guiones de trabajo forman parte del repositorio.

Los archivos de resultados no se versionan junto con el código. El repositorio contiene los guiones que los producen y los verificadores que los leen, pero no sus salidas, y tampoco el corpus. La razón es que el material de partida son archivos cifrados por ransomware real y notas de rescate auténticas, cuya publicación no corresponde. Los resultados quedan disponibles como material suplementario de acceso restringido, del que se excluyen los archivos que contienen indicadores de compromiso. El corpus de notas cuenta además con un manifiesto de trazabilidad que documenta, para cada nota, su fuente, su tipo de procedencia y su extensión original.
